\documentclass[11pt]{article}
\usepackage[utf8]{inputenc}
\usepackage{amsmath, amssymb, amsthm}
\usepackage[margin=1in]{geometry}

\theoremstyle{definition}
\newtheorem{definition}{Definition}
\newtheorem{theorem}{Theorem}
\newtheorem{lemma}{Lemma}

\theoremstyle{remark}
\newtheorem{remark}{Remark}

\DeclareMathOperator{\tr}{tr}
\newcommand{\dual}[2]{{}_{V'}\langle #1, #2 \rangle_{V}}

\title{Lecture 8: Gelfand Triples and Abstract Parabolic Equations}
\date{04.12.2025}
\author{Tien Anh Vu}
\begin{document}
\maketitle

\section*{Chapter 3: Variational approach to SPDEs (continued)}

\subsection*{Gelfand triples}

Let $V$ and $H$ be separable Hilbert spaces with a dense and continuous embedding
\[
V \hookrightarrow H.
\]
Concretely, ``continuous embedding'' means that there is a constant $C>0$ such that
\[
\|u\|_H \le C\|u\|_V,
\qquad
u\in V,
\]
so the $V$-norm is stronger than the $H$-norm (more regularity in $V$ implies control in $H$).
Without loss of generality, we may assume that
\[
\|u\|_H \le \|u\|_V,
\qquad
u\in V.
\]
This normalization is harmless: if originally $\|u\|_H\le C\|u\|_V$, one can replace the norm on $V$
by the equivalent norm $\|u\|_{V,\mathrm{new}}:=C\|u\|_V$.
Equivalently, one may extend the $V$-norm to all of $H$ by setting $\|u\|_V:=\infty$ for
$u\in H\setminus V$; then the inequality above holds for every $u\in H$.
By identifying $H$ with its dual $H'$ via the Riesz isomorphism, we obtain the \emph{Gelfand triple}
\[
V \hookrightarrow H \cong H' \hookrightarrow V'.
\]
The meaning of the last embedding $H'\hookrightarrow V'$ is that every functional on $H$ restricts
to a functional on $V$; this restriction is continuous because $V\hookrightarrow H$ is continuous.
We write $\dual{\cdot}{\cdot}$ for the duality pairing between $V'$ and $V$; it extends the inner product on $H$ in the sense that
\[
\dual{h}{v} = \langle h,v\rangle_H,
\qquad
h\in H,\ v\in V.
\]

\begin{remark}[Typical example]
For a bounded domain $D\subset\mathbb{R}^d$,
\[
V = H_0^1(D),
\qquad
H = L^2(D),
\qquad
V' = H^{-1}(D).
\]
In this example, $V$ is the energy space, $H$ is the pivot space with the $L^2$ inner product,
and $V'$ is the corresponding dual space; concretely, $H_0^1(D)$ is the Sobolev space of
$L^2$-functions with weak first derivatives in $L^2$ and zero boundary values, while
$H^{-1}(D)$ is its dual, the negative Sobolev space.
\end{remark}

\begin{remark}[Why one says ``space regularity'']
In PDE applications, the space $V$ usually describes regularity in the spatial variable. For
instance, if $V=H_0^1(D)$, then $u(t)\in V$ means that for each fixed time $t$ the function
$x\mapsto u(t,x)$ has one weak spatial derivative in $L^2(D)$, so membership in
$L^2([0,T];V)$ expresses spatial regularity integrated over time.
\end{remark}

\subsection*{Operators $A:V\to V'$}

Let $A:V\to V'$ be a (typically linear) operator. Common assumptions in the variational theory are:
\begin{itemize}
  \item \textbf{Boundedness:} $A\in \mathcal{L}(V,V')$, i.e.\ there exists $C>0$ such that
  \[
  \|Au\|_{V'}\le C\|u\|_V,
  \qquad
  u\in V.
  \]
  A lot of the time, the operators arising in PDEs are unbounded when viewed as operators on $H$.
  In the variational framework, the relevant condition is that $A$ is bounded as a map from $V$
  into $V'$.
  Intuitively, boundedness means that $A$ does not amplify the size of $u$ in an uncontrolled
  way: if $u$ is small in $V$, then $Au$ is small in $V'$.
  \item \textbf{Coercivity (bounded from below):} there exist $\alpha>0$ and $\lambda\in\mathbb{R}$ such that
  \[
  \dual{Au}{u}\ge \alpha\|u\|_V^2-\lambda\|u\|_H^2,
  \qquad
  u\in V.
  \]
  Intuitively, coercivity says that $A$ controls the stronger $V$-norm of $u$ up to a lower-order
  correction involving the weaker $H$-norm; this is the key estimate behind energy bounds.
  Depending on the sign convention in the evolution equation, one often also encounters the equivalent
  ``dissipative'' form
  \[
  \dual{Au}{u}\le -\alpha\|u\|_V^2+\lambda\|u\|_H^2,
  \]
  which is obtained by replacing $A$ with $-A$ (e.g.\ writing the heat equation as $u_t-\Delta u=f$ rather than $u_t=\Delta u+f$).
  \item \textbf{Monotonicity:} for all $u,v\in V$,
  \[
  \dual{Au-Av}{u-v}\ge 0.
  \]
  Intuitively, monotonicity means that $A$ does not push two states farther apart in a harmful
  direction; it is a nonlinear analogue of positivity and is central for uniqueness and stability.
\end{itemize}

\subsection*{Example: the Laplacian}

Take $V=H_0^1(D)$ and define $A:V\to V'$ by
\[
\dual{Au}{v}
:=
\int_D \nabla u\cdot \nabla v\,dx,
\qquad
u,v\in H_0^1(D).
\]
Formally, this corresponds to $Au=-\Delta u$ with homogeneous Dirichlet boundary conditions, since for smooth $u,v$,
\[
\int_D \nabla u\cdot \nabla v\,dx
=
-\int_D (\Delta u)\,v\,dx.
\]
Moreover, this definition shows directly that $A$ is bounded from $V$ to $V'$: for $u\in V$,
\[
\|Au\|_{V'}
=
\sup_{\|v\|_V\le 1}\dual{Au}{v}
=
\sup_{\|v\|_V\le 1}\int_D \nabla u\cdot \nabla v\,dx
\le
\sup_{\|v\|_V\le 1}\|\nabla u\|_{L^2(D)}\|\nabla v\|_{L^2(D)}
\le
\|u\|_V.
\]
Here we used Cauchy--Schwarz and the fact that $\|v\|_V$ controls $\|\nabla v\|_{L^2(D)}$ for $v\in H_0^1(D)$.

\subsection*{Abstract parabolic Cauchy problem}

\begin{theorem}[Existence and uniqueness in a Gelfand triple]
Let $A\in\mathcal{L}(V,V')$ be coercive, let $u_0\in H$, and let $f\in L^2([0,T];V')$. Then the evolution problem
\[
\begin{cases}
u'(t)=Au(t)+f(t), & t\in(0,T),\\
u(0)=u_0,
\end{cases}
\]
has a unique solution
\[
u\in C([0,T];H)\cap L^2([0,T];V),
\qquad
u'\in L^2([0,T];V').
\]
In words, $u$ has space-time regularity: it is continuous in time with values in $H$ and
square-integrable in time with values in the more regular space $V$, while its time derivative is
square-integrable with values in the dual space $V'$.
\end{theorem}

\begin{remark}[Heat equation]
With $V=H_0^1(D)$, $H=L^2(D)$, and $A=-\Delta$, the abstract problem reads
\[
u'(t)=\Delta u(t)+f(t),
\qquad
u(0)=u_0,
\]
and the solution satisfies $u(t)\in H_0^1(D)$ for a.e.\ $t$ and $u\in C([0,T];L^2(D))$.
This means that for almost every fixed time $t$, the spatial function $x\mapsto u(t,x)$ belongs to
$H_0^1(D)$, i.e.\ it has weak first derivatives in $L^2(D)$ and zero boundary values.
\end{remark}

\begin{remark}[Why this is called parabolic]
If $A=-\Delta$, then the abstract equation is the heat equation
\[
u_t-\Delta u=f.
\]
In one space dimension this becomes $u_t-u_{xx}=f$, whose second-order part has discriminant
$A=-1$, $B=0$, $C=0$, hence $B^2-AC=0$, so it is parabolic in the classical classification.
In higher dimensions one says it is parabolic because it is first order in time and its spatial
second-order part $-\Delta$ is elliptic.
\end{remark}

\subsection*{Uniqueness: energy estimate}

To prove uniqueness, let $u$ and $v$ be two solutions in the class
\[
C([0,T];H)\cap L^2([0,T];V),
\qquad
u',v'\in L^2([0,T];V'),
\]
with the same initial condition $u(0)=v(0)$. Set $w:=u-v$. Then $w$ solves the corresponding homogeneous
parabolic equation
\[
\begin{cases}
w'(t)=Aw(t), & t\in(0,T),\\
w(0)=0.
\end{cases}
\]
The key identity is the following.

\begin{lemma}[Derivative of the $H$-norm]
If $w\in L^2([0,T];V)$ and $w'\in L^2([0,T];V')$, then the map $t\mapsto \|w(t)\|_H^2$ is absolutely continuous and for a.e.\ $t\in(0,T)$,
\[
\frac{d}{dt}\|w(t)\|_H^2
=
2\dual{w'(t)}{w(t)}.
\]
\end{lemma}

Applying this with $w'=Aw$ yields, for a.e.\ $t$,
\[
\frac{d}{dt}\|w(t)\|_H^2
=
2\dual{Aw(t)}{w(t)}.
\]
If, for example, $A$ satisfies the dissipative estimate
\[
\dual{Aw(t)}{w(t)}\le \lambda \|w(t)\|_H^2,
\]
Here $\lambda\in\mathbb{R}$ is a constant measuring how much growth is still allowed in the
$H$-energy. Intuitively, this one-sided estimate says that $A$ does not increase the size of $w$
too fast in $H$. This is different from coercivity, which gives a lower bound on
$\dual{Au}{u}$ and controls the stronger $V$-norm of $u$.
then
\[
\frac{d}{dt}\|w(t)\|_H^2 \le 2\lambda \|w(t)\|_H^2.
\]
Equivalently,
\[
\frac{d}{dt}\Bigl(e^{-2\lambda t}\|w(t)\|_H^2\Bigr)\le 0.
\]
Hence the map $t\mapsto e^{-2\lambda t}\|w(t)\|_H^2$ is nonincreasing. Since $w(0)=0$, we obtain
\[
\|w(t)\|_H^2 = 0
\qquad \text{for all } t\in[0,T].
\]
Hence $w\equiv 0$, and therefore $u=v$.

\subsection*{Existence via approximation (Galerkin method)}

Let $(e_k)_{k\ge 1}$ be a complete orthonormal system (CONS: $C$ means
$\overline{\mathrm{span}\{e_1,e_2,\dots\}}=H$, $O$ means
$\langle e_i,e_j\rangle_H=\delta_{ij}$, and $S$ means that $(e_k)_{k\ge 1}$ is a family
of elements of $H$) of $H$ such that $e_k\in V$ for all $k$.
For $n\in\mathbb{N}$ define the finite-dimensional subspace
\[
V_n := \mathrm{span}\{e_1,\dots,e_n\}\subset V.
\]
We project the equation onto $V_n$ and look for an approximate solution $u_n\in C([0,T];V_n)$ of
the form
\[
u_n(t)=\sum_{k=1}^n g_k^{(n)}(t)\,e_k,
\]
where the scalar functions $g_k^{(n)}$ are the \emph{Fourier coefficients} of $u_n$ in the basis $(e_k)$.
Then $u_n$ is the unique solution in $C([0,T];V_n)$ of an $n$-dimensional linear ordinary
differential equation. More precisely, testing the evolution equation against $e_k$,
$k=1,\dots,n$, yields the finite-dimensional system
\[
\begin{cases}
\dfrac{d}{dt}\langle u_n(t),e_k\rangle_H = \dual{Au_n(t)}{e_k} + \dual{f(t)}{e_k}, & t\in(0,T),\\[6pt]
\langle u_n(0),e_k\rangle_H = \langle u_0,e_k\rangle_H, &
\end{cases}
\qquad k=1,\dots,n.
\]
This is an $n$-dimensional Cauchy problem for the coefficient vector
\[
\bigl(g_1^{(n)},\dots,g_n^{(n)}\bigr),
\]
so by the standard existence and uniqueness theory for linear ODEs it has a unique solution.

\begin{remark}[On choosing the basis $(e_k)$]
Since $V$ is dense in $H$, one can choose a CONS $(e_k)$ with $e_k\in V$ (e.g.\ by applying Gram--Schmidt to a dense sequence in $V$).
In PDE examples it is common to take $(e_k)$ as eigenfunctions of an elliptic operator (such as the Dirichlet Laplacian), which makes the Galerkin system coincide with a truncated Fourier expansion.
\end{remark}

\begin{remark}[What elliptic means here]
For a second-order linear operator
\[
Lu=-\sum_{i,j=1}^d a_{ij}(x)\,\partial_{ij}u+\text{lower-order terms},
\]
elliptic means that the coefficient matrix $A(x)=(a_{ij}(x))$ is positive definite, i.e.
\[
\sum_{i,j=1}^d a_{ij}(x)\xi_i\xi_j>0
\qquad\text{for every }\xi\in\mathbb{R}^d\setminus\{0\}.
\]
For the Laplacian, $a_{ij}(x)=\delta_{ij}$, so $-\Delta$ is elliptic.
\end{remark}

\subsection*{Uniform energy estimate (a priori bound)}

It is called an \emph{a priori} estimate because the bound is derived from the equation and the
data alone, before knowing the solution explicitly.

The next step is to pass to the limit $n\to\infty$. The essential ingredient for this is an a
priori estimate that is uniform in $n$.
Testing the Galerkin equation with $u_n(t)\in V_n$ and using coercivity yields an energy inequality of the form
\[
\sup_{n\ge 1}\left(\sup_{t\in[0,T]}\|u_n(t)\|_H^2 \;+\; \int_0^T \|u_n(t)\|_V^2\,dt\right) < \infty,
\]
In words, all approximate solutions $u_n$ satisfy the same bound: their $H$-norms stay bounded
for all times $t\in[0,T]$, and their $V$-norms are square-integrable in time. The constant in
this estimate depends only on $\|u_0\|_H$, $\|f\|_{L^2([0,T];V')}$, and the constants in the
assumptions on $A$, but not on $n$.

\begin{remark}[Heat equation: why Sobolev spaces appear]
For the heat equation on a domain $D\subset\mathbb{R}^d$ with $V=H_0^1(D)$ and $H=L^2(D)$, the $V$-norm is (up to equivalence) the $L^2$-norm of the gradient:
\[
\|u\|_V^2 \simeq \|\nabla u\|_{L^2(D)}^2.
\]
Here $\simeq$ means that the two quantities are comparable up to multiplicative constants: each
one bounds the other. This follows from Poincar\'e's inequality on $H_0^1(D)$.

\begin{remark}[Why this follows from Poincar\'e's inequality]
If we take
\[
\|u\|_V^2=\|u\|_{H_0^1(D)}^2:=\|u\|_{L^2(D)}^2+\|\nabla u\|_{L^2(D)}^2,
\]
then Poincar\'e's inequality gives
\[
\|u\|_{L^2(D)} \le C_P\|\nabla u\|_{L^2(D)},
\qquad u\in H_0^1(D).
\]
Therefore
\[
\|u\|_V^2
=
\|u\|_{L^2(D)}^2+\|\nabla u\|_{L^2(D)}^2
\le
(C_P^2+1)\|\nabla u\|_{L^2(D)}^2.
\]
Conversely,
\[
\|\nabla u\|_{L^2(D)}^2 \le \|u\|_V^2.
\]
This is immediate from the definition of $\|\cdot\|_V$, since $\|u\|_V^2$ is the sum of
$\|u\|_{L^2(D)}^2$ and $\|\nabla u\|_{L^2(D)}^2$.
Hence
\[
\|\nabla u\|_{L^2(D)}^2 \le \|u\|_V^2 \le (C_P^2+1)\|\nabla u\|_{L^2(D)}^2,
\]
which is exactly the meaning of
\[
\|u\|_V^2 \simeq \|\nabla u\|_{L^2(D)}^2.
\]
\end{remark}

For the heat equation, this gives a concrete example of the uniform a priori estimate: for the
Galerkin sequence $(u_n)$,
\[
\sup_{n\ge 1}\left(\sup_{t\in[0,T]}\|u_n(t)\|_{L^2(D)}^2+\int_0^T\|\nabla u_n(t)\|_{L^2(D)}^2\,dt\right)<\infty.
\]
This is why one works in Sobolev spaces: for $u\in H_0^1(D)$, the gradient $\nabla u$ is defined
in the weak sense, and the estimate gives control of these weak derivatives.
\end{remark}

\subsubsection*{Energy estimate: derivation}

We now prove the uniform a priori estimate for the Galerkin approximations $(u_n)$.

Starting from the Galerkin system, one computes (formally; justified since $u_n$ is finite-dimensional and smooth in time)
Since $(e_1,\dots,e_n)$ is an orthonormal basis of $V_n$ and $u_n(t)\in V_n$, Parseval's identity gives
\[
\|u_n(t)\|_H^2=\sum_{k=1}^n \langle u_n(t),e_k\rangle_H^2.
\]
Therefore
\[
\frac12\frac{d}{dt}\|u_n(t)\|_H^2
=
\frac12\frac{d}{dt}\sum_{k=1}^n \langle u_n(t),e_k\rangle_H^2
=
\sum_{k=1}^n \langle u_n'(t),e_k\rangle_H\,\langle u_n(t),e_k\rangle_H.
\]
Using $u_n'(t)=Au_n(t)+f(t)$ and the defining property of the Galerkin approximation,
\[
\sum_{k=1}^n \langle u_n'(t),e_k\rangle_H\,\langle u_n(t),e_k\rangle_H
=
\sum_{k=1}^n \dual{Au_n(t)+f(t)}{e_k}\,\langle u_n(t),e_k\rangle_H
\]
Since
\[
u_n(t)=\sum_{k=1}^n \langle u_n(t),e_k\rangle_H\,e_k,
\]
linearity of the duality pairing in the second argument gives
\[
\sum_{k=1}^n \dual{Au_n(t)+f(t)}{e_k}\,\langle u_n(t),e_k\rangle_H
=
\dual{Au_n(t)+f(t)}{\sum_{k=1}^n \langle u_n(t),e_k\rangle_H\,e_k}
=
\dual{Au_n(t)+f(t)}{u_n(t)}.
\]
Hence
\[
\frac12\frac{d}{dt}\|u_n(t)\|_H^2
=
\dual{Au_n(t)}{u_n(t)}
\;+\;
\dual{f(t)}{u_n(t)}.
\]
Assuming the simplified coercivity $\dual{Au}{u}\le -\alpha\|u\|_V^2$ (or after rewriting $A$ by shifting a lower-order term), and applying Cauchy--Schwarz in $(V',V)$ followed by Young's inequality,
\[
\dual{f(t)}{u_n(t)}
\le
\|f(t)\|_{V'}\|u_n(t)\|_V
\le
\frac{1}{2\alpha}\|f(t)\|_{V'}^2+\frac{\alpha}{2}\|u_n(t)\|_V^2,
\]
where we used Young's inequality in the general form
\[
ab\le \frac{a^p}{p}+\frac{b^q}{q},
\qquad a,b\ge 0,\quad \frac1p+\frac1q=1,
\]
or, in the weighted quadratic case relevant here,
\[
ab\le \frac{1}{2\varepsilon}a^2+\frac{\varepsilon}{2}b^2.
\]
Here we take $p=q=2$, $a=\|f(t)\|_{V'}$, $b=\|u_n(t)\|_V$, and $\varepsilon=\alpha$; equivalently,
one may view this as the weighted AM--GM inequality.
Therefore
\[
\frac12\frac{d}{dt}\|u_n(t)\|_H^2
\le
-\alpha\|u_n(t)\|_V^2
\;+\;
\frac{1}{2\alpha}\|f(t)\|_{V'}^2
\;+
\frac{\alpha}{2}\|u_n(t)\|_V^2
=
-\frac{\alpha}{2}\|u_n(t)\|_V^2
\;+\;
\frac{1}{2\alpha}\|f(t)\|_{V'}^2.
\]
Integrating from $0$ to $t$ gives
\[
\frac12\|u_n(t)\|_H^2-\frac12\|u_n(0)\|_H^2
\le
-\frac{\alpha}{2}\int_0^t\|u_n(s)\|_V^2\,ds
\;+\;
\frac{1}{2\alpha}\int_0^t\|f(s)\|_{V'}^2\,ds.
\]
Equivalently,
\[
\|u_n(t)\|_H^2
\;+\;
\alpha\int_0^t\|u_n(s)\|_V^2\,ds
\le
\|u_n(0)\|_H^2
\;+\;
\frac{1}{\alpha}\int_0^t\|f(s)\|_{V'}^2\,ds.
\]
Replacing $t$ by the supremum over $[0,T]$ and weakening the constants, we obtain the a priori estimate
\[
\sup_{t\in[0,T]}\|u_n(t)\|_H^2
\;+\;
\frac{\alpha}{2}\int_0^T\|u_n(t)\|_V^2\,dt
\le
\|u_n(0)\|_H^2 + \frac{1}{2\alpha}\int_0^T\|f(t)\|_{V'}^2\,dt.
\]
Since $u_n(0)$ is the $H$-orthogonal projection of $u_0$ onto $V_n$, $\|u_n(0)\|_H\le \|u_0\|_H$, and the right-hand side is independent of $n$.
Using this uniform bound and weak compactness, we may extract a subsequence, again denoted by
$(u_n)$, such that
\[
u_n\rightharpoonup u \quad \text{in } L^2([0,T];V).
\]
If $A\in\mathcal{L}(V,V')$ is linear, then also
\[
A u_n \rightharpoonup A u \quad \text{in } L^2([0,T];V').
\]
Indeed, let $A^*:V\to V'$ denote the adjoint operator, characterized by
\[
\dual{Aw}{v}=\dual{A^*v}{w}
\qquad\text{for all }w,v\in V.
\]
Then for any $v\in L^2([0,T];V)$,
\[
\int_0^T \dual{Au_n(t)}{v(t)}\,dt
=
\int_0^T \dual{A^*v(t)}{u_n(t)}\,dt.
\]
Since $A^*\in\mathcal{L}(V,V')$, we have $A^*v\in L^2([0,T];V')$, and by weak convergence of
$u_n$ in $L^2([0,T];V)$,
\[
\int_0^T \dual{A^*v(t)}{u_n(t)}\,dt
\longrightarrow
\int_0^T \dual{A^*v(t)}{u(t)}\,dt
\qquad\text{for all }v\in L^2([0,T];V).
\]
Applying the adjoint identity once more, this limit equals
\[
\int_0^T \dual{Au(t)}{v(t)}\,dt.
\]
This allows one to pass to the limit in the variational formulation and obtain a solution of the
abstract parabolic problem.

\subsection*{Generalization to nonlinear operators}

Assume now that $A:V\to V'$ is still bounded and coercive but \emph{not necessarily linear}.
The same a priori estimate remains true, so after extracting a subsequence we may assume
\[
u_n \rightharpoonup u \quad \text{in } L^2([0,T];V).
\]
Moreover, the growth bound on $A$ implies that $(A(u_n))$ is bounded in $L^2([0,T];V')$, so,
after passing to a further subsequence if necessary,
\[
A(u_n) \rightharpoonup g \quad \text{in } L^2([0,T];V')
\]
for some (a priori unknown) $g$.
In contrast to the linear case, weak convergence $u_n\rightharpoonup u$ does \emph{not} directly imply
$A(u_n)\rightharpoonup A(u)$, so we need an additional argument to identify $g$.

\subsection*{Monotonicity (dissipativity)}

Assume the following one-sided Lipschitz/monotonicity condition: there exists
$\lambda\in\mathbb{R}$ such that for all $u,v\in V$,
\[
\dual{A(u)-A(v)}{u-v}
\le
\lambda\|u-v\|_H^2.
\]
Together with this, assume the growth bound
\[
\|A(u)\|_{V'}\le M(1+\|u\|_V),
\qquad u\in V,
\]
and the hemicontinuity condition: for all $u,v,w\in V$, the map
$\theta\mapsto \dual{A(u+\theta v)}{w}$ is continuous on $\mathbb{R}$.

For any $v\in L^2([0,T];V)$, the one-sided Lipschitz/monotonicity condition gives, after
integrating in time,
\[
\int_0^T \dual{A(u_n(t))-A(v(t))}{u_n(t)-v(t)}\,dt
\le
\lambda\int_0^T \|u_n(t)-v(t)\|_H^2\,dt.
\]
Expanding this inequality gives
\[
\begin{aligned}
&\int_0^T \dual{A(u_n(t))}{u_n(t)}\,dt
\;-\;
\int_0^T \dual{A(u_n(t))}{v(t)}\,dt\\
&\quad
\;-\;
\int_0^T \dual{A(v(t))}{u_n(t)}\,dt
\;+\;
\int_0^T \dual{A(v(t))}{v(t)}\,dt\\
&\le \lambda\int_0^T \|u_n(t)-v(t)\|_H^2\,dt.
\end{aligned}
\]
From the energy estimate, one obtains the $\limsup$ bound
\[
\limsup_{n\to\infty}\int_0^T \dual{A(u_n(t))}{u_n(t)}\,dt
\le
\int_0^T \dual{g(t)}{u(t)}\,dt.
\]
Also,
\[
\int_0^T \dual{A(u_n(t))}{v(t)}\,dt
\longrightarrow
\int_0^T \dual{g(t)}{v(t)}\,dt
\]
by weak convergence of $A(u_n)$ in $L^2([0,T];V')$, and
\[
\int_0^T \dual{A(v(t))}{u_n(t)}\,dt
\longrightarrow
\int_0^T \dual{A(v(t))}{u(t)}\,dt
\]
because $A(v)\in L^2([0,T];V')$ and $u_n\rightharpoonup u$ in $L^2([0,T];V)$. Passing to the
limit yields
\[
\int_0^T \dual{g(t)-A(v(t))}{u(t)-v(t)}\,dt
\le
\lambda\int_0^T \|u(t)-v(t)\|_H^2\,dt
\qquad
\forall\,v\in L^2([0,T];V).
\]

The remaining step is the classical \emph{Browder--Minty (Minty) trick}: we perturb $v$ around $u$
and let the perturbation parameter tend to zero, using hemicontinuity to pass to the limit in the
nonlinear term.

Now choose
\[
v=u+\theta w,
\qquad \theta>0,\quad w\in L^2([0,T];V).
\]
Then $u-v=-\theta w$ and therefore
\[
-\theta\int_0^T \dual{g(t)-A(u(t)+\theta w(t))}{w(t)}\,dt
\le
\lambda\theta^2\int_0^T \|w(t)\|_H^2\,dt.
\]
Dividing by $\theta>0$ and multiplying by $-1$ yields
\[
\int_0^T \dual{g(t)-A(u(t)+\theta w(t))}{w(t)}\,dt
\ge
-\lambda\theta\int_0^T \|w(t)\|_H^2\,dt.
\]
Letting $\theta\downarrow 0$ and using hemicontinuity gives
\[
\int_0^T \dual{g(t)-A(u(t))}{w(t)}\,dt \ge 0
\qquad
\forall\,w\in L^2([0,T];V).
\]
Replacing $w$ by $-w$ gives
\[
\int_0^T \dual{g(t)-A(u(t))}{-w(t)}\,dt \ge 0,
\qquad \forall\,w\in L^2([0,T];V),
\]
equivalently,
\[
\int_0^T \dual{g(t)-A(u(t))}{w(t)}\,dt \le 0,
\qquad \forall\,w\in L^2([0,T];V).
\]
Therefore,
\[
\int_0^T \dual{g(t)-A(u(t))}{w(t)}\,dt = 0
\qquad
\forall\,w\in L^2([0,T];V),
\]
hence $g=A(u)$ in $L^2([0,T];V')$.

\medskip
We now relate the quantity $\int_0^T \dual{A(u_n(t))}{u_n(t)}\,dt$ to the energy balance and use
it to compare $\int_0^T \dual{g(t)}{u(t)}\,dt$ with $\liminf_{n\to\infty}\int_0^T \dual{A(u_n(t))}{u_n(t)}\,dt$.

First, the Galerkin equation can be written as $u_n'(t)=A(u_n(t))+f(t)$ in $V_n$.
Testing it by $u_n(t)$ and using $\frac12\frac{d}{dt}\|u_n(t)\|_H^2=\dual{u_n'(t)}{u_n(t)}$ gives
\[
\dual{A(u_n(t))}{u_n(t)}
=
\frac12\frac{d}{dt}\|u_n(t)\|_H^2
\;-\;
\dual{f(t)}{u_n(t)}.
\]
Integrating over $[0,T]$ yields
\[
\int_0^T \dual{A(u_n(t))}{u_n(t)}\,dt
=
\frac12\bigl(\|u_n(T)\|_H^2-\|u_n(0)\|_H^2\bigr)
\;-\;
\int_0^T \dual{f(t)}{u_n(t)}\,dt.
\]
\medskip
Next, passing to the limit in the Galerkin formulation yields (in $V'$)
\[
u'(t)=g(t)+f(t),
\]
meaning that for a.e.\ $t\in[0,T]$ and all $\varphi\in V$,
\[
\dual{u'(t)}{\varphi}=\dual{g(t)}{\varphi}+\dual{f(t)}{\varphi}.
\]
Choosing $\varphi=u(t)$, integrating over $[0,T]$, and using
$\frac12\frac{d}{dt}\|u(t)\|_H^2=\dual{u'(t)}{u(t)}$, we obtain
\[
\int_0^T \dual{g(t)}{u(t)}\,dt
=
\frac12\bigl(\|u(T)\|_H^2-\|u(0)\|_H^2\bigr)
\;-\;
\int_0^T \dual{f(t)}{u(t)}\,dt.
\]
\medskip
Finally, extract a subsequence such that $u_n(T)\rightharpoonup u(T)$ in $H$.
Since $u_n\rightharpoonup u$ in $L^2([0,T];V)$, we have
\[
\int_0^T \dual{f(t)}{u_n(t)}\,dt \longrightarrow \int_0^T \dual{f(t)}{u(t)}\,dt,
\]
and since $u_n(0)\to u(0)$ in $H$ (Galerkin projection of the common initial datum) together with
weak lower semicontinuity of the norm in $H$, we infer
\[
\begin{aligned}
\int_0^T \dual{g(t)}{u(t)}\,dt
&\le
\liminf_{n\to\infty}\Biggl[
\frac12\bigl(\|u_n(T)\|_H^2-\|u_n(0)\|_H^2\bigr)
\;-\;
\int_0^T \dual{f(t)}{u_n(t)}\,dt
\Biggr]\\
&=
\liminf_{n\to\infty}\int_0^T \dual{A(u_n(t))}{u_n(t)}\,dt\\
&\le
\limsup_{n\to\infty}\int_0^T \dual{A(u_n(t))}{u_n(t)}\,dt.
\end{aligned}
\]
Thus we have linked the limit quantity $\int_0^T\dual{g(t)}{u(t)}\,dt$ to the $\liminf/\limsup$ of
$\int_0^T \dual{A(u_n(t))}{u_n(t)}\,dt$.
Combining this with the monotonicity inequality and passing to the limit in $n$ yields, for every
$v\in L^2([0,T];V)$,
\[
\int_0^T \dual{g(t)-A(u(t))}{u(t)-v(t)}\,dt
\le
\limsup_{n\to\infty}\int_0^T \dual{A(u_n(t))-A(v(t))}{u(t)-v(t)}\,dt
\le
0.
\]
This variational inequality is the key output of the Browder--Minty/Minty trick and is the step
that leads to the identification $g=A(u)$.

\subsection*{Compactness}

The monotonicity method above is not the only way to identify nonlinear terms in a limit.
Another important approach is to use \emph{compactness} to upgrade weak convergence to strong convergence in an intermediate space.

Suppose that $(u_n)$ is bounded in
\[
L^2([0,T];V).
\]
In practice one also needs a bound that prevents fast oscillations in time. A common formulation is a fractional time-regularity estimate such as
\[
u_n \text{ bounded in } L^2([0,T];V)\cap H^{1/2}([0,T];V'),
\]
or, in many PDE/SPDE settings, a bound on the time derivative $u_n'$ in a negative space (e.g.\ $u_n'$ bounded in $L^2([0,T];V')$).
If the embedding $V\hookrightarrow H$ is compact, then (under additional time-regularity bounds, e.g.\ on $(u_n')$ in a negative space) one can often extract a subsequence that converges \emph{strongly} in
\[
L^2([0,T];H).
\]
Statements of this type are instances of the Aubin--Lions/Simon compactness lemma: compactness in space plus some regularity in time yields compactness in space--time.
Recall that the embedding $V\hookrightarrow H$ is called \emph{compact} if the inclusion map
$i:V\to H$, $i(v)=v$, is a compact linear operator, i.e.\ it maps bounded sets in $V$ to
\emph{relatively compact} sets in $H$.
Here ``relatively compact in $H$'' means that every sequence in the set admits a subsequence that
converges in $H$ (equivalently, the closure of the set is compact in $H$).

\begin{remark}[PDE example: Rellich--Kondrachov]
For a bounded domain $D\subset\mathbb{R}^d$,
\[
H_0^1(D)\hookrightarrow L^2(D)
\]
is a compact embedding (Rellich--Kondrachov theorem): any sequence bounded in $H_0^1(D)$ admits a
subsequence converging strongly in $L^2(D)$.
This is a prototypical source of compactness in variational PDE/SPDE theory.
\end{remark}

\begin{remark}[Why compactness helps]
Strong convergence in $L^2([0,T];H)$ is especially useful when the equation contains nonlinearities that are not monotone (e.g.\ convection terms in fluid dynamics such as Navier--Stokes).
In such cases, weak convergence may be insufficient to pass to the limit in products, while strong convergence can restore convergence of nonlinear terms.
\end{remark}

\end{document}
